\scanpage{625}
\mspage{616}

% L1
Now let
\[
\tag{\unsure{20}}
\mu_{\tau'}=\{1,\tau'\}\subset\mathcal{G}^{\wedge}
\]
% L3
then \inserted{\unsure{$\mu_{\tau'}$ normalizes $L_\rho$}},
\deleted{\unsure{this}} for
\(\tau'(\rho)=\rho^{-1}\) whence \(\tau'(L_\rho)=L_\rho\)
hence it normalizes \(Z_\rho\unsure{\equiv|=}\operatorname{Cent}_{\mathcal{M}}(L_\rho)\), hence
also \(Z_\rho^!=Z_\rho\cap\mathcal{M}^!\)
(since \(\mathcal{M}^!\) is invariant \unsure{in} \(\mathcal{M}\))
\(=Z_{\rho,0}^!\times\mu=\mathcal{M}'(-\bar{\jmath})\times\mu\).
% L9
\deleted{Let} Moreover
\(\mu_{\tau'}\cap Z_\rho^!\subset
\deleted{\unsure{\{1\}}}\mu_{\tau'}\cap\mathcal{M}^!=\{1\}\).
Let
\[
\tag{21}
\mathcal{N}_\rho^?=\mu_{\tau'}\cdot Z_\rho^!
=\mu_{\tau'}\cdot
\underbrace{(Z_{\rho,0}^!\times\mu)}_{
\substack{\mathcal{M}(-\bar{\jmath})\times\mu\\
\simeq\Gamma_{\mathbb{Q}}^{\prime\sim}\times\mu}}
\quad\deleted{\unsure{\simeq\mathcal{M}_{\mathbb{Q}}\times\mu|\simeq\mathcal{M}_{0}\times\mu}}.
\]
\hfill ($\frac12$ direct product)

% L13
One should be aware that contrary to what
this notation might suggest,
\(\mu_{\tau'}\) does \uline{not} normalize \(Z_{\rho,0}^!\) --- otherwise one
would have a \unsure{decomposition|arrangement} into a product
\[
\mathcal{N}_\rho^?=\mathcal{N}_{\rho,0}^?\times\mu,
\quad\text{where }\mathcal{N}_{\rho,0}^?
=\deleted{\unsure{Z_{\rho,0}^!\times\mu}}
\mu_{\tau'}\cdot\unsure{Z_{\rho,0}^?|Z_{\rho,0}^!},
\]
% L19
but this is not the case. We have an exact sequence
\[
\tag{22}
\diagram{\begin{tikzcd}[ampersand replacement=\&,column sep=large,row sep=small]
1\arrow[r] \& \mu\arrow[r] \& \mathcal{N}_\rho^?\arrow[r]
\& \gamma_{\rho,\sigma}\arrow[r,no head]\arrow[d,equal] \& 1\\
\& \& \& \Gamma_{\mathbb{Q}}^{\sim} \&
\end{tikzcd}}
\]
% L22
so that \(\mathcal{N}_\rho^?\) is a quasi-section
of multiplicity $2$ of the extension
\(\mathcal{N}_\rho^{*}\longrightarrow\gamma_{\rho,\sigma}\), --- of
\(\mathcal{M}\longrightarrow\gamma_{\rho,\sigma}=\Gamma_{\mathbb{Q}}^{\sim}\).
Thus we find

\margin{624}

\scanpage{626}

% L1
\[
\tag{23}
\left\{\begin{aligned}
\mathcal{M}&=\mathcal{N}_\rho^?\wedge_{\mu_2}\mathcal{G}^{+\wedge}\\
\unsure{\mathcal{N}_\rho^*}&=\mathcal{N}_\rho^?\wedge_{\mu_2}
\deleted{\unsure{\mathbb{Z}}}L_\rho
\simeq\mathcal{N}_\rho^?\times
\underbrace{L_{-\rho}}_{\simeq\mathbb{Z}/3\mathbb{Z}}
\end{aligned}\right.
\]
% L3
The \deleted{sequence}\inserted{extension} (22) of
\(\gamma_{\rho,\sigma}\) by \(\mu\) is
the extension denoted by \(\widetilde{\gamma}_{\rho\sigma}\) in the preceding section,
it is surely \uline{\unsure{not} \deleted{trivial}\inserted{\unsure{split}}}.
It gives rise to a canonical \unsure{homomorphism}
\[
\tag{24}
\diagram{\begin{tikzcd}[ampersand replacement=\&,column sep=large]
\gamma_{\rho,\sigma}\arrow[r,"\sim"]
\&\mathcal{N}_\rho^?/\mu\arrow[r]
\&\mathcal{N}_\rho^*/\mu.
\end{tikzcd}}
\]

% L9
\uline{Remark} To do it properly, one would have to
make explicit the operation of \(\tau'\) on
\(Z_\rho^!\simeq Z_{\rho,0}^!\times\mu
\xrightarrow{\sim}\Gamma_{\mathbb{Q}}^{\prime\sim}\times\mu\), in terms of
the operation of the image \(\tau_{\Gamma\deleted{\unsure{\mathbb Q}}}\) of \(\tau'\) on
\(\Gamma_{\mathbb Q}^{\prime\sim}\) by inner \unsure{automorphism}
(N.B. \(\tau_\Gamma\) is none other than the element of
\(\Gamma_{\mathbb Q}\subset\Gamma_{\mathbb Q}^{\sim}\), ``complex conjugation'')
and of a \unsure{homomorphism} from \(\Gamma_{\mathbb Q}^{\prime\sim}\) to \(\mu\)
\inserted{i.e. of a ``quadratic character'' of \(\Gamma_{\mathbb Q}^{\prime\sim}\)}
--- it is a good bet that
this is precisely \(\varepsilon_2\) --- so that
\(\Gamma_{\mathbb Q}^{0\sim}\simeq
\Gamma_{\mathbb Q}^{\prime\sim}\cap\Gamma_{\mathbb Q}^{\prime\prime\sim}\)
appears as invariant under \(\tau'\) \ldots
In short, it would be well to take up again the study
of the extension made in the preceding section,
by \unsure{starting there} from the point of view of the \unsure{associated} \(\alpha,\beta\),

\margin{625}

\scanpage{627}
\mspage{617}

% L1
in the present perspective, probably
simpler. It would be well at the same time
to \deleted{define}\inserted{make explicit} another partial section,
which \deleted{would deserve} would be over
\(\deleted{\unsure{\Gamma}}\Gamma_{\mathbb Q}^{\prime\prime\prime\sim}\),
\[
\Gamma_{\mathbb Q}^{\prime\prime\prime\sim}
=\{x\in\Gamma_{\mathbb Q}^{\sim}\mid
\varepsilon_2(x)\varepsilon_3(x)=1\},
\]
% L6
and which would perhaps deserve the notation
\(\mathcal{M}(\sqrt3)\) \unsure{or at least}
\(\mathcal{M}(q)\), for
a suitable generator \(\unsure{q|\rho}\) of \(\mathbb{Q}(\sqrt3)\) ???
\margin{?\quad\uline{\unsure{Ponder these}}
\uline{\unsure{coincidences.} \unsure{Attentive.|attentively}}}

% L9
I realize that I am
\unsure{talking crap} --- having somewhat lost touch
with the geometric content of my
calculations. One must keep in mind that
\[
\tag{25}
\diagram{\begin{tikzcd}[ampersand replacement=\&]
\mathcal{M}_0^!\arrow[r,"\sim"] \& \mathcal{M}_{0,3}^{!\sim}
\end{tikzcd}}
\]
% L14
corresponds \unsure{up to} $\sim$, to the \(\pi_1\) of \(U_{0,3\mathbb Q}\),
the sections of \(\mathcal{M}_0^!\longrightarrow
\Gamma_{\mathbb Q}^{\sim}\simeq\gamma_{\rho,\sigma}\) thus correspond
(up to $\sim$) to those of
\(\pi_1(U_{0,3\mathbb Q})\) over \(\Gamma_{\mathbb Q}\),
thus (``\unsure{invariably|inevitably|unavowably}'') to the
\inserted{\unsure{rational/$\mathbb Q$}} points of \(U_{0,3\mathbb Q}\)
--- when they are partial sections, to the
points of \(U_{0,3\mathbb Q}\) with values in finite
extensions of \(\mathbb Q\). Thus, \deleted{the} the section
\deleted{\unsure{privileged}}\inserted{considered}
in b) does indeed correspond to the \(\mathrm{pt}\ \frac12\) of \(U_{0,3\mathbb Q}\),
the \deleted{\unsure{quasi-}}section
\inserted{partial \(\mathcal M(-\bar\jmath)\)} considered here in c) corresponds to the
point \(-\bar\jmath\) of \(U_{0,3\mathbb Q}\), which is in
\(U_{0,3\deleted{\unsure{\mathbb Q}}}(\mathbb Q[\sqrt[3]{1}])\),
so the notation \inserted{\(\mathcal M(-\bar\jmath)\)} is reasonable.
As for \(\mathcal N_\rho^{\sim}\),
it is \uline{not} \(\subset\mathcal M_0^!\) --- even on dividing by \(\mu\);
one finds after division by \(\mu\) a section

\margin{626}

\scanpage{628}

% L1
of \(\mathcal M_{0,3}^{\sim}\) over \(\Gamma_{\mathbb Q}^{\sim}\), but this is not a
section of \(\mathcal M_{0,3}^{!\sim}\). Now up to $\sim$, \(\mathcal M_{0,3}^{\sim}\)
is the \(\pi_1\) of \(M_{1,1\mathbb Q}'\), moduli scheme over \(\mathbb Q\)
of elliptic curves \uline{mod symmetry}.
Thus giving this section corresponds to
giving such an ``elliptic curve
mod symmetry'' over \(\mathbb Q\). Passing to the \deleted{sub}
subgroup \(\Gamma_{\mathbb Q}^{\prime\sim}\) --- which corresponds to the extension
of the base field from \(\mathbb Q\) to \(\mathbb Q(j)=\mathbb Q(\sqrt[3]{1})\)
one recovers the image in \(\mathcal M_{0,3}^{!\sim}\) of the preceding \deleted{quasi-}
section \inserted{partial} --- \deleted{it} to \(\mathcal M(-\bar\jmath)\)
--- this thus shows that it is still the ``same'' elliptic
curve. I thus have the impression \deleted{that} however
that over the whole of \(\Gamma_{\mathbb Q}^{\sim}\), this section of
\(\mathcal M_{0,3}^{\sim}\) cannot be lifted to a section
of \(\unsure{\mathcal M_\mu|\mathcal M}\) over \(\Gamma_{\mathbb Q}^{\sim}\)
--- i.e. precisely that the extension \(\mathcal N_\rho^{\unsure{?|*}}\)
of \(\Gamma_{\mathbb Q}^{\sim}\) by \(\mu\) is nontrivial.
Would this \uline{mean} that there is no elliptic curve
defined/\(\mathbb Q\) with invariant $0$? (Something that seems absurd
--- one must be able over \uline{a field} to give any
invariant\ldots). \margin{?}
To be precise, we have the
diagram of moduli schemes over \(\mathbb Q\)
\[
\tag{26}
\diagram{\begin{tikzcd}[ampersand replacement=\&,column sep=large,row sep=large]
\begin{gathered}\text{moduli scheme}\\\text{of elliptic}\\\text{curves}\end{gathered}
\&M_{1,1}\arrow[d,"\substack{\text{gerbe}\\\text{of group }\mu}"']\arrow[dr]
\& \&\\
\begin{gathered}\text{moduli scheme of}\\\text{elliptic curves}\\\text{mod symmetry}\end{gathered}
\&M_{1,1}'\arrow[r]
\&\underset{\unsure{0|\circ}}{\mathbb E^1}
\&\begin{gathered}\text{coarse moduli}\\\text{scheme}\end{gathered}
\end{tikzcd}},
\]
% L26
the fact that a point of \(M_{1,1}'(\mathbb Q)\) does not li-

\margin{627}

\scanpage{629}
\mspage{618}

% L1
ft to a point of \(M_{1,1}\), does not \uline{mean}
that its image in \(\mathbb E^1\) does not lift
to \(M_{1,1}\). To be precise, the \deleted{\illegible{1}}\inserted{diagram}
identifies with:
\[
\tag{27}
\diagram{\begin{tikzcd}[ampersand replacement=\&,column sep=large,row sep=large]
M_{1,1}\simeq(U_{0,3},\widetilde{\mathcal G}_3)
  \arrow[d]\arrow[dr]\&\\
M_{1,1}'\simeq(U_{0,3},\mathcal G_3)\arrow[r]
  \&U_{0,3}/\mathcal G_3\simeq U_{0,3}/\widetilde{\mathcal G}_3
     \simeq\mathbb E^1
\end{tikzcd}},
\]
% L7
where \(\widetilde{\mathcal G}_3\) is the extension of
\(\mathcal G_3\)\inserted{\(=\mathrm{SL}(2,\mathbb Z/2\mathbb Z)\)}
by \(\mu\), defined by
\[
\tag{28}
\begin{aligned}
\widetilde{\mathcal G}_3
 &=\mathcal G^+/\pi_0
 \simeq\mathrm{SL}(2,\mathbb Z)/
 \begin{gathered}
 \text{invariant subgroup}\\
 \text{generated by}\\[-2pt]
 -\ell_0
 \end{gathered}
 \\
 &\simeq\mathrm{SL}(2,\mathbb Z/4\mathbb Z)/
 \begin{gathered}
 \text{invariant}\\
 \text{subgroup}\\
 \text{(of order 4) generated}\\
 \text{by }-\ell_0,\ \text{i.e. \unsure{subgroup|subgroup}}\\
 \text{generated by}\\
 -\ell_0,-\ell_1.
 \end{gathered}
\end{aligned}
\]
% L13
by letting \(\widetilde{\mathcal G}_3\) operate on \(U_{0,3}\) through
\(\mathcal G_3\). A point of \(M_{1,1}'\)\inserted{\unsure{rational/\(\mathbb Q\)}}
identifies with \(\mathcal G_3\)-torsor \inserted{\(T\)} over \(\mathbb Q\), equipped
with a map
\[
\tag{29}
\diagram{\begin{tikzcd}[ampersand replacement=\&]T\arrow[r]\&U_{0,3}\end{tikzcd}}
\]
\margin{\(T\) is the torsor
\inserted{\unsure{under \(\mathrm{SL}(2,\mathbb Z/2\mathbb Z)\)}}
of \unsure{rigidifications|privilegiations} of the points of order $2$ and
\(\widetilde T\) that of ``\unsure{pri-rigidifications|privilegiations|rigidifications}
of the points of order \(\unsure{\leq4|4}\)
\unsure{made precise}'' (the first is defined
for an \unsure{elliptic} curve mod symmetry, and
the $2^{\text{nd}}$ only for a genuine elliptic
curve)}
% L18
compatible with the actions of \(\mathcal G_3\), a
point of \deleted{\(U_{0,3}\)} \(M_{1,1}\) identifies with a \(\widetilde{\mathcal G}_3\)-torsor
\(\widetilde T\) over \(\mathbb Q\), equipped with a map
\[
\tag{30}
\diagram{\begin{tikzcd}[ampersand replacement=\&]\widetilde T\arrow[r,no head]\&U_{0,3}\end{tikzcd}}
\]
compatible with the operation of \(\widetilde{\mathcal G}_3\). \uline{Denoting}
by \deleted{\unsure{\(T\)}}
\[
T=\widetilde T/\mu
\]
the corresponding \(\mathcal G_3\)-torsor, giving
(30) is equivalent to giving (29). Thus the

\margin{628}

\scanpage{630}

% L1
liftings to \(M_{1,1}\) of a point of \(M_{1,1}'\) rational/\(\mathbb Q\)
identify with the \inserted{\(\widetilde{\mathcal G}_3\)-}torsors \(\widetilde T\)
which lift the \(\mathcal G_3\)-torsor \(T\), i.e. (up to isomorphism) they correspond
to the homomorphisms \(\Gamma_{\mathbb Q}\longrightarrow\widetilde{\mathcal G}_3\)
which lift the homomorphism
\(\Gamma_{\mathbb Q}\longrightarrow\mathcal G_3\) which defines \(T\):
\margin{we suppose \unsure{now} chosen
a \deleted{\unsure{geometric}} point of \(T(\overline{\mathbb Q}_0)\)}
\[
\diagram{\begin{tikzcd}[ampersand replacement=\&,column sep=large,row sep=small]
\&\widetilde{\mathcal G}_3\arrow[d]\\
\Gamma_{\mathbb Q}\arrow[r]\arrow[ur,dashed]\&\mathcal G_3
\end{tikzcd}}
\quad -
\]
% L8
the obstruction is, as it should be, an
extension of \(\Gamma_{\mathbb Q}\) by \(\mu_2\), i.e. an \unsure{element}
of \(H^2(\mathbb Q,\mu_2)\simeq {}_2\mathrm{Br}(\mathbb Q)\).
In the case at hand, the image of \(T\) in \(U_{0,3}\)
is made up of \(\{-\bar\jmath,\unsure{-j|-\bar\jmath}\}\)
which is connected over \(\mathbb Q\),
\uline{hence \(\Gamma_{\mathbb Q}\longrightarrow\mathcal G_3\) is surjective}.
\margin{\unsure{corrected|erroneous}, cf. correction below}
In general, \deleted{\unsure{one must}} in the Galois-theoretic translation
in terms of splittings
of \(\mathcal M_{0,3}\) over \(\Gamma_{\mathbb Q}\), the homomorphism
\(\Gamma_{\mathbb Q}\longrightarrow\mathcal G_3\)
is the composite
\[
\diagram{\begin{tikzcd}[ampersand replacement=\&]
\Gamma_{\mathbb Q}\arrow[r,"\unsure{k|\kappa}"]
\&\mathcal M_{0,3}\arrow[r]\&\mathcal G_3
\end{tikzcd}},
\]
% L22
which is trivial iff \deleted{\unsure{\(\Gamma_{\mathbb Q}\)}}
\inserted{\(\unsure{k|\kappa}\)} is a section of
\deleted{of}
\[
\mathcal M_{0,3}^!=\operatorname{Ker}
 (\mathcal M_{0,3}\longrightarrow\mathcal G_3)
\]
over \(\Gamma_{\mathbb Q}\). Here we have a section of
\(\mathcal M_{0,3}^!\) \deleted{over}\inserted{above}
the subgroup \(\Gamma_{\mathbb Q}'\) of index $2$,
\deleted{hence the image of \(\Gamma_{\mathbb Q}'\)}
\(\unsure{k|\kappa}\mid\Gamma_{\mathbb Q}'=1\), hence
\(\unsure{k|\kappa}(\Gamma_{\mathbb Q})\subset\mathcal G_3\)
is of order $2$, we see that

\margin{629}

\scanpage{631}
\mspage{619}

it is the subgroup \inserted{$\{1,\dot{\sigma}_{\infty}\}$} of $\mathcal{G}_3$.
(for the image of $\tau'\inserted{{}=\sigma\tau}\in\mathcal{M}_{0,3}$ in $\mathcal{G}_3$
\deleted{\unsure{is}} (since that of $\tau$ is $1$) is equal to that of
$\sigma=\sigma_{\infty}$, i.e. $\dot{\sigma}_{\infty}$. This means
that \deleted{\unsure{T}} we have
\[
T\simeq T_0\mathbin{\wedge}_{\{1,\dot{\sigma}_{\infty}\}}\mathcal{G}_3
\]
i.e. $T$ is obtained from a $\mu$-torsor (i.e. from a
quadratic extension of $\mathbb{Q}$\unsure{)} --- namely precisely
$\mathbb{Q}(\sqrt[3]{1})=\mathbb{Q}(j))$ by extension of the groups
of operators \inserted{$\mu\longrightarrow\mathcal{G}_3\quad(-1\longmapsto\dot{\sigma}_{\infty})$}.
Thus, taking the inverse image of $\mu=\{1,\dot{\sigma}_{\infty}\}$ in
$\widetilde{\mathcal{G}}_3$, we find an extension of $\mu$ by $\mu$
--- it can hardly be anything but $\mathbb{Z}/4\mathbb{Z}$ ! --- which I feel like denoting $\mathcal{E}''$

\[
\diagram{\begin{tikzcd}[ampersand replacement=\&,column sep=1.3em,row sep=0.15em]
\&\& \mathbb{Z}/4\mathbb{Z}\arrow[d,equal] \&\&\\
1\arrow[r]\&\mu\arrow[r]\&\mathcal{E}''\arrow[r]\&\mu\arrow[r,no head]\&1
\end{tikzcd}}
\]

\margin{yes, this is obvious since $\sigma=\sigma_{\infty}$ in
$\widetilde{\mathcal{G}}_3=\mathrm{SL}(2,\mathbb{Z})/\unsure{\Pi_0|\Pi}$
is indeed of order $4$, so $\mathcal{E}''\simeq\mathbb{Z}/4\mathbb{Z}$
\unsure{canonically}\ldots}
and the \unsure{problem} of obstruction is to find a lifting

\[
\diagram{\begin{tikzcd}[ampersand replacement=\&,column sep=3em,row sep=2em]
\mathcal{E}''\arrow[r,no head]\&\mu\\
\Gamma_{\mathbb Q}\arrow[u,dashed]\arrow[ur]\&
\end{tikzcd}}
\]
,
(which \deleted{\unsure{we}} surely does not exist) --- so
the obstruction is identified as the extension of $\Gamma_{\mathbb Q}$ by $\mu$,
inverse image of the extension $\mathcal{E}''$.
\margin{This is obvious in \unsure{effect|fact}, \unsure{that we|since} :
$(\Gamma_{\mathbb Q})_{\mathrm{ab}}\simeq\widehat{\mathbb Z}^{*}
\supset\mathbb Z_3^{*}$, now
$(\mathbb Z_3^{*})_4\simeq\mathbb Z/2\mathbb Z$
(\unsure{so already on} $\mathbb Z_3^{*}$ \unsure{it does not lift}\ldots}
(One should look at why, in

\margin{630}

\scanpage{632}

classical terms, this obstruction is \unsure{split}, when
one passes to the subgroup $\Gamma_{\mathbb Q}^{\prime\prime\prime}$ corresponding to the
field $\mathbb Q(\sqrt{3}))$

Thus, the $\mathbb Q$-rational point under consideration of
$M_{1,1}'$ does not lift to a $\mathbb Q$-rational point of
$M_{1,1}$ --- On the other hand, its image $\inserted{0}$ in
$\mathbb E^1\simeq U_{0,3}/\mathcal G_3$ must indeed
lift --- i.e. there must indeed exist a
$\widetilde{\mathcal G}_3$-torsor $\widetilde P$, and a
$\widetilde{\mathcal G}_3$-morphism
\deleted{$\unsure{P}$}

\[
\diagram{\begin{tikzcd}[ampersand replacement=\&,column sep=1.6em,row sep=0.1em]
\widetilde P\arrow[r]\&U_{0,3}\&\text{i.e.}\&P\arrow[r]\arrow[d,equal]\&U_{0,3}\\
\&\&\&\widetilde P/\mu\&
\end{tikzcd}}
\]
.

Let us consider indeed the \inserted{sub-}scheme
\deleted{of $U_{0,3}$, inverse image}
\[
\deleted{\unsure{Q=P}}\qquad
Q=U_{0,3}\mathbin{\times}_{\mathbb E^1}\{0\}
\]
the inverse image of $0$, i.e.
\[
P_0=Q_{\mathrm{red}}\ \deleted{\unsure{\simeq}}
\]
it is a subscheme which, over a field
$K\supset\mathbb Q(j)$, \deleted{splits into} is a sum of
two points, and which over $\mathbb Q$ is connected
--- it is a scheme \inserted{canonical\textsuperscript{ly}} isomorphic to
$\operatorname{Spec}\mathbb Q(-\bar j)$. It is stable under $\mathcal G_3$,
which operates on it via

\[
\diagram{\begin{tikzcd}[ampersand replacement=\&]
\mathcal G_3\arrow[r,"\mathrm{sg}"]\&\mu
\end{tikzcd}}
\]

and the Galois operation of $\mu$ on the $\mu_{\mathbb Q}$-torsor $P_0$.

\margin{631}

\scanpage{633}
\mspage{620}

Thus we have the diagram

\[
\diagram{\begin{tikzcd}[ampersand replacement=\&,column sep=1.8em,row sep=2em]
\&\widetilde{\mathcal G}_3\arrow[r,"\mathrm{can}"]\&\mathcal G_3\arrow[r,"\mathrm{sg}"]\&\mu\\
(31)\&\&\Gamma_{\mathbb Q}\arrow[ur]\&
\end{tikzcd}}
\]

\margin{\uline{N.B.} the kernel of
$\widetilde{\mathcal G}_3\longrightarrow\mu$ is a
\inserted{central} extension of $\mathbb Z/3\mathbb Z$
by $\mu=\mathbb Z/2\mathbb Z$, it
is therefore \unsure{trivial} and
isomorphic to $\mathbb Z/6\mathbb Z$
(kernel generated by $\rho$).}

We thus see that the category of liftings
of the $\mathbb Q$-rational point $0$ of $\mathbb E^1$ to a
$\mathbb Q$-rational point of $M_{1,1}'$ resp. $M_{1,1}$, is
equivalent to that of $\mathcal G_3$-torsors \inserted{$P$} resp\unsure{.}
of $\widetilde{\mathcal G}_3$-torsors \inserted{$\widetilde P$} over $\mathbb Q$, which lift the
$\mu$-torsor $P_0$. We had chosen a
$P=T$ in a \unsure{particularly} \unsure{trivial|trivial} way,
using the splitting

\[
\diagram{\begin{tikzcd}[ampersand replacement=\&,column sep=1.5em]
\mu\arrow[r,"\sim"]\&\{1,\dot\sigma_{\infty}\}\arrow[r,hook]\&\mathcal G_3
\end{tikzcd}}
\]

of $\mathcal G_3\longrightarrow\mu$ --- but the torsor $T$ does not
lift to a $\widetilde T$. The isomorphism classes of points of $M_{1,1}'$ resp. of $M_{1,1}$ over
the point $0$ of $\mathbb E^1$, \unsure{correspond} likewise to the
conjugacy classes \inserted{(modulo kernels \inserted{$\mathcal G_3^+,\widetilde{\mathcal G}_3^+$} of $\mathcal G_3\longrightarrow\mu$ and $\widetilde{\mathcal G}_3\longrightarrow\mu$)} of liftings of
$\Gamma_{\mathbb Q}\longrightarrow\mu$ to $\Gamma_{\mathbb Q}\longrightarrow\mathcal G_3$ resp. to $\Gamma_{\mathbb Q}\longrightarrow\widetilde{\mathcal G}_3$.

\deleted{\unsure{From there}} Let us consider a lifting to

\[
\diagram{\begin{tikzcd}[ampersand replacement=\&]
\Gamma_{\mathbb Q}\arrow[r]\&\mathcal G_3
\end{tikzcd}}
\]
,
its image is a subgroup of $\mathcal G_3$, which is either

\margin{632}

\scanpage{634}

$\{1,\dot\sigma_{\infty}\}$ or one of its conjugates,
$\{1,\dot\sigma_0\}$, $\{1,\unsure{\bar\sigma_1|\dot\sigma_1}\}$,
\uline{or} it is the whole of $\mathcal G_3$. In the first
case, we have seen that it does not lift
to $\widetilde{\mathcal G}_3$. So to be able to lift,
one must take an \uline{epimorphism} $\Gamma_{\mathbb Q}\longrightarrow\mathcal G_3$ ---
so we leave the abelian setting over
$\mathbb Q$, to enter the domain of
abelian extensions of $\mathbb Q(j)$.

To examine in terms of Galois
groups whether there exists such a lifting \deleted{which}
$\Gamma_{\mathbb Q}\longrightarrow\widetilde{\mathcal G}_3$, and \deleted{\unsure{which}}
how to obtain them,
appears as a pleasant
and \unsure{delectable} exercise in class field
theory, which I am not going to pursue
for the moment, so as not to make
this digression last forever. Returning later
to this point, I would need at the same
time to make explicit the relation between the
\unsure{choices|points} of a lifting of the homomorphism from $\Gamma_{\mathbb Q}$,
and splittings of the extension $\mathcal M_{0,3}$ of $\Gamma_{\mathbb Q}$
by $\unsure{\mathcal G_{0,3}^{+\wedge}|\widetilde{\mathcal G}_{0,3}^{+\wedge}}$,
or $\mathcal M$ of $\Gamma_{\mathbb Q}$ by
$\mathcal G^{+\wedge}\simeq\mathrm{SL}(2,\mathbb Z)^\wedge$.

\margin{633}

\scanpage{635}
\mspage{621}

Let us return to the subgroups \deleted{$\unsure{\mathcal N_{\rho}^{?}}$ of}

\[
\diagram{
\begin{tikzcd}[ampersand replacement=\&,column sep=3em,row sep=2em]
\mathcal N_{\rho}^{?}
  \arrow[r,phantom,"{\overset{3}{\unsure{\subset}}}"{description}]
\&\mathcal N_{\rho}^{*}\arrow[r,phantom,"{\subset}"{description}]
\&\mathcal M\\
Z_{\rho,0}^{!}=\mathcal M'(-\bar j)
  \arrow[u,phantom,"{4\,\cup}"{description}]
\&\&
\end{tikzcd}
}
\]

containing \unsure{the|the} partial \deleted{section}\inserted{\unsure{subgroup|subgroups}}
$\mathcal M(-\bar j)\simeq\widetilde\Gamma_{\mathbb Q}'$.

Let
\[
\tag{32}
\begin{aligned}
\mathcal M^0(-\bar j)&\overset{2}{\subset}\mathcal M'(-\bar j)\\
\underset{\unsure{\mathrm{dfn}|\mathrm{def}}}{=}&\operatorname{Ker}
\bigl(\mathcal M'(-\bar j)\xrightarrow{\varepsilon_2}\mu\bigr)
=\deleted{\unsure{\mathcal N_{\rho}^{?}|\mathcal M}}\,Z_{\rho}
\cap\deleted{\unsure{\operatorname{Ker}\varepsilon_2}}\,
\mathcal M_0^{\prime\prime!}\\
&=\mathcal N_{\rho}^{*}\cap\mathcal M_0^{0!}
\end{aligned}
\]
which is therefore a partial section \unsure{up|up} above
\[
\tag{33}
\begin{aligned}
\widetilde\Gamma_{\mathbb Q}^0
&=\operatorname{Ker}\bigl(\widetilde\Gamma_{\mathbb Q}
\xrightarrow{(\varepsilon_2,\varepsilon_3)}\mu\times\mu\bigr)
\deleted{\unsure{\ne|=}\widetilde\Gamma_{\mathbb Q}}\\
&=\operatorname{Ker}\bigl(\widetilde\Gamma_{\mathbb Q}
\longrightarrow\widehat{\mathbb Z}^{*}\longrightarrow
(\mathbb Z/12\mathbb Z)^{*}\bigr)
\end{aligned}
\]

\deleted{Let us consider} Modulo a verification
which I am going to make below, $\tau'$ normalizes
\deleted{(34)}
$\mathcal M^0(-\bar j)$, let us consider \unsure{the|the}
\[
\tag{34}
\mathcal M'''(-\bar j)\deleted{\unsure{\mathcal N_{\rho,0}^{?}}}
\overset{\unsure{\mathrm{dfn}|\mathrm{def}}}{=}
\underset{\text{(semidirect)}}{\mu_{\tau'}\mathbin{.}\mathcal M^0(-\bar j)}
\qquad\text{where }\mu_{\tau'}=\{1,\tau'\}
\]
it is a subgroup of index $4$ of $\mathcal N_{\rho}^{?}$,
and we now have
\[
\tag{35}
\deleted{\unsure{\mathcal N_{\rho,0}^{?}}}\quad
\mathcal M'''(-\bar j)\xrightarrow{\sim}\widetilde\Gamma_{\mathbb Q}''',
\]
it is a partial section \unsure{up|up} above
$\widetilde\Gamma_{\mathbb Q}'''$, whence isomorphisms

\margin{634}

\scanpage{636}

\deleted{
\[
\tag{36}
\left\{
\begin{aligned}
\mathcal M'''&\simeq\mathcal M'''(-\bar j)\mathbin{.}\mathcal G^{+\wedge}\\
\deleted{\unsure{\mathcal M'''}}\,\mathcal N_{\rho}^{*}
&\simeq\mathcal M'''(-\bar j)\mathbin{.}L_{\rho}.
\end{aligned}
\right.
\]

We have the summary diagram
of inclusions of subgroups of $\mathcal N_{\rho}^{*}$

\[
\diagram{
\begin{tikzcd}[ampersand replacement=\&,column sep=2.5em,row sep=1.5em]
\&\&\&Z_{\rho}\arrow[dr,hook,"2"]\&\\
\&\&\begin{gathered}\mathcal M'(-\bar j)\\=Z_{\rho,0}^{!}\end{gathered}
  \arrow[ur,hook,"6"]\arrow[dr,hook,"4"]
\&\&\mathcal N_{\rho}^{*}\\
(37)\&\mathcal M^0(-\bar j)\arrow[ur,hook,"2"]\arrow[dr,hook,"2"']
\&\&\mathcal N_{\rho}^{?}\arrow[ur,hook,"3"']\&\\
\&\&\mathcal M'''(-\bar j)\arrow[ur,hook,"4"]\&\&
\end{tikzcd}
}
\]


Let us note that the inclusion
\[
\tag{38}
\underset{\substack{\simeq\\\unsure{\mathbb D_6|\mathbb D_{\sigma}}}}{\mathbb D_{\rho}}
\overset{\mathrm{dfn}}{=}
\underset{\simeq\mathbb Z/2\mathbb Z}{\uline{\mu_{\tau'}}}
\mathbin{.}
\underset{=\mathbb Z/6\mathbb Z}{\uline{L_{\rho}}}
\subset\mathcal N_{\rho}^{*}
\]
induces an isomorphism of
\[
\tag{39}
\mathbb D_{\rho}\xrightarrow{\sim}\mathcal N_{\rho}^{*}/
\]
}

\margin{635}

\scanpage{637}
\mspage{622}

\[
\tag{36}
\left\{
\begin{aligned}
\mathcal{M}'''&\simeq \mathcal{M}'''(-\bar j)\mathbin{.}\mathcal{G}^{+\wedge}\\
\mathcal{N}_{\rho}^{*}&\simeq \mathcal{M}'''(-\bar j)\mathbin{.}L_{\rho}.
\end{aligned}
\right.
\]

Summary diagram of the subgroups of $\mathcal{N}_{\rho}^{*}$
obtained so far:


\[
\diagram{
\begin{tikzcd}[ampersand replacement=\&,column sep=1.05em,row sep=2.8em,column sep=9pt,row sep=14pt,execute at begin picture={\fontsize{10}{12}\selectfont}]
\&\& \mathcal{M}^{0}(-\bar j)
  \arrow[rr,hook,"2"]
  \arrow[dl,hook,"2"']
\&\& \begin{gathered}\mathcal{M}'(-\bar j)\\=Z_{\rho,0}^{!}\end{gathered}
  \arrow[rr,hook,"2"]
  \arrow[dl,hook,"2"']
  \arrow[dr,no head,draw=gray,
    "{\unsure{\mathcal{M}'(-\bar j)\times\mu}}"{description,pos=.5},
    "{\unsure{\searrow}}"{pos=.93,xshift=10pt,yshift=-6pt}]
\&\& \mathcal{M}'(-\bar j)\times\mu
  \arrow[rr,hook,"3"]
  \arrow[dl,hook,"2"']
\&\& Z_{\rho}\arrow[dl,hook,"2"']\\
(37)\&\mathcal{M}'''(-\bar j)\arrow[rr,hook,"2"]
\&\& \mathcal{M}'''(-\bar j)\times\mu\arrow[rr,hook,"2"{pos=.75}]
\&\& \mathcal{N}_{\rho}^{?}\arrow[rr,hook,"3"{pos=.75}]
\&\& \mathcal{N}_{\rho}^{*}
\&\&\\
\&\&\&\&1
  \arrow[uu,hook]
  \arrow[uull,hook,start anchor=north west,end anchor=south west,bend left=25]
  \arrow[rr,hook,"2"]
  \arrow[dl,hook,"2"']
\&\&\mu\arrow[uu,hook]\arrow[rr,hook,"3"]\arrow[dl,hook,"2"']
\&\&L_{\rho}\arrow[uu,hook]\arrow[dl,hook,"2"']\\
\&\&\&\mu_{\tau'}\arrow[uu,hook]\arrow[uull,hook]\arrow[rr,hook,"2"]
\&\&\mathbb{D}_{\tau'}\arrow[uu,hook]\arrow[rr,hook,"3"]
\&\&\mathbb{D}_{\rho}\arrow[uu,hook]
\end{tikzcd}
}
\]
,

where we are here interested \unsure{above all} in the
``upper face'', \unsure{containing|concerning} subgroups of finite index in
$\mathcal{N}_{\rho}^{*}$, intermediate between it and
$\unsure{\mathcal{M}_{0}(-\bar j)|\mathcal{M}^{0}(-\bar j)}$
(of index $24$ in $\mathcal{N}_{\rho}^{\unsure{*}}$).
I am going to examine the question of invariance of these subgroups
in one Another, and the structure of the quotient groups.

I am going to redo a diagram which \unsure{really|better} shows the
\deleted{\unsure{groups|relationships}} relations between the groups in question.

\margin{636}

\scanpage{638}


\[
\diagram{
\begin{tikzcd}[ampersand replacement=\&,column sep=.6em,row sep=1.3em,cells={font=\scriptsize},nodes in empty cells]
\&\&Z_{\rho}\arrow[rrr,hook,"2"]
  \arrow[ddddll,dashed,leftrightarrow,bend right=35,
    "{\unsure{\leftrightarrow}}"'{description}]
\&\&\&\mathcal{N}_{\rho}^{*}\\
\&\&\&Z_{\rho,0}^{!}=\mathcal{M}'(-\bar j)
  \arrow[rrr,hook,"2"]
  \arrow[dl,hook,"2"']
\&\&\&\mathcal{M}'(-\bar j)\mathbin{.}\mu_{\tau'}\;{}^{*}
  \arrow[dl,hook,"2"']\\
\&\&Z_{\rho}^{!}=\mathcal{M}'(-\bar j)\times\mu
  \arrow[uu,hook,"3"']
  \arrow[rrr,hook,"2"]
  \arrow[ddll,hook,dashed,"3"']
\&\&\&\mathcal{N}_{\rho}^{?}
  \arrow[uu,hook,"3"'{pos=.7}]
  \arrow[ddll,hook,dashed,"3"']\\
\&\&\&\mathcal{M}^{0}(-\bar j)
  \arrow[uu,hook,"2"'{pos=.7}]
  \arrow[rrr,hook,"2"]
  \arrow[ddl,hook,"2"']
\&\&\&\mathcal{M}'''(-\bar j)
  \arrow[uu,hook,"2"']
  \arrow[ddl,hook,"2"']\\
Z_{\rho}\arrow[rrr,hook,dashed]
\&\&\&\unsure{\mathcal{N}_{\rho}|\mathcal{N}_{\rho}^{*}}
  \arrow[uuuurr,dashed,leftrightarrow,out=0,in=-36,out distance=78pt,in distance=59pt,start anchor=east,end anchor=south east,
    "{\unsure{\leftrightarrow}}"'{description,pos=.25}]\\
\&\&\mathcal{M}^{0}(-\bar j)\times\mu
  \arrow[uuu,hook,"2"']
  \arrow[rrr,hook,"2"]
  \arrow[dddll,hook,
    "{\deleted{\unsure{\hookrightarrow}}}"{description}]
\&\&\&\mathcal{M}'''(-\bar j)\times\mu
  \arrow[uuu,hook,"2"']
  \arrow[dddll,hook,
    "{\deleted{\unsure{\hookrightarrow}}}"{description}]\\
\&\&\&1\arrow[uuu,hook,bend right=50]
  \arrow[rrr,hook,"2"]
  \arrow[dl,hook,"2"']
\&\&\&\mu_{\tau'}\arrow[uuu,hook]\arrow[dl,hook,"2"']\\
(38)\&\&\mu\arrow[uu,hook]\arrow[rrr,hook,"2"]
  \arrow[ddll,hook,dashed,"3"']
\&\&\&\mathbb{D}_{\tau'}\arrow[uu,hook]
  \arrow[ddll,hook,dashed,"3"']\\
\deleted{\unsure{L_{\rho}}}
  \arrow[rrr,hook,
    "{\deleted{\unsure{\hookrightarrow}}}"{description}]
\&\&\&\deleted{\unsure{\mathbb{D}_{\rho}}}\\
L_{\rho}\arrow[uuuuu,hook,dashed,bend left=20]
  \arrow[rrr,hook,dashed,"2"]
\&\&\&\mathbb{D}_{\rho}\arrow[uuuuu,hook,dashed,bend right=20]
\end{tikzcd}
}
\]


\margin{$*$ -- Caution,
$\mathcal{M}'(-\bar j)\mathbin{.}\mu_{\tau'}$ is \uline{not}
a subgroup, for $\mu_{\tau'}$ does not normalize
$\mathcal{M}'(-\bar j)$.}

\unsure{a)} Subgroups of finite index in $Z_{\rho}$ -- in the diagram
there are four of them


\[
\diagram{
\begin{tikzcd}[ampersand replacement=\&,column sep=3em,row sep=3em]
\&Z_{\rho}^{!}=\mathcal{M}'(-\bar j)\times\mu
\&\mathcal{M}'(-\bar j)=Z_{\rho,0}^{!}=\mathcal{N}_{\rho,0}^{\unsure{*!}}
  \arrow[l,hook,"2"']\\
(39)\&\mathcal{M}^{0}(-\bar j)\times\mu
  \arrow[u,hook,"2"']
\&\mathcal{M}^{0}(-\bar j)=Z_{\rho,0}^{0!}
  \arrow[u,hook,"2"]
  \arrow[l,hook,"2"]
\end{tikzcd}
}
\]


are they invariant in $Z_{\rho}$? This is clear for the subgroups
\deleted{$\unsure{\mathcal{M}'(-\bar j)}$}
\[
\tag{40}
\left\{
\begin{aligned}
\mathcal{M}'(-\bar j)\times\mu
  &=Z_{\rho}\cap\mathcal{M}^{!}=Z_{\rho}^{!},
  \quad\text{and for}\\
\deleted{\unsure{\mathcal{M}^{0}(-\bar j)\times\mu=}}&\\[-.5ex]
\mathcal{M}^{0}(-\bar j)\times\mu
  &=Z_{\rho}\cap\mathcal{M}^{!}\cap
       \operatorname{Ker}(\varepsilon_{2}:\mathcal{M}\to\mu)
\end{aligned}
\right.
\]
-- this shows that these groups are

\margin{637}

\scanpage{639}
\mspage{623}

even invariant in $\mathcal{N}_{\rho}^{*}$, since
\deleted{$\unsure{\mathcal{M}^{!}}$}.
$Z_{\rho},\mathcal{M}^{!},\operatorname{Ker}\varepsilon_{2}$
are stabilized by $\mathcal{N}_{\rho}^{*}$.
Besides, \deleted{since}
$Z_{\rho,0}^{!}=\mathcal{M}'(-\bar j)$ is
\deleted{of index $2$}\inserted{invariant}
\deleted{\unsure{in}}\inserted{in}
$Z_{\rho}^{!}=\mathcal{M}'(-\bar j)\times\mu$,
\unsure{to} show that it is so in $Z_{\rho}$, it suffices to prove
that it is normalized by $\rho$
(since $Z_{\rho}=L_{\rho}\mathbin{.}Z_{\rho}^{!}$ -- this is
$\pm$ trivial), now let us examine in general for
\deleted{$\unsure{x\in}$} $x\in\mathcal{M}$
whether $x$ normalizes \deleted{$\unsure{\mathcal{M}_{0}^{!}}$}
\[
\mathcal{M}_{0}^{!}\overset{\mathrm{dfn}}{=}
\mathcal{M}(0)\mathbin{.}\Pi_{0},
\]
Since it normalizes $\Pi_{0}$, it suffices to see whether
for $u\in\mathcal{M}(0)$, we have
\[
xux^{-1}\in\mathcal{M}_{0}^{!}
\]
which is also equivalent to
\[
xux^{-1}\equiv u\qquad(\Pi_{0})
\]
i.e.
\[
[x,u]\in\Pi_{0}
\]
now we know that for $x\in\mathcal{G}$,
\inserted{$u\in\mathcal{M}^{!}$,} we have
\[
\tag{*}
\underbrace{[x,u]}_{\inserted{x\,u(x)^{-1}}}^{\unsure{c}}
\equiv
(\operatorname{sg}(x))^{\inserted{\unsure{\xi_{2}(u)|\varepsilon_{2}(u)}}}
\deleted{\unsure{\xi_{2}(u)|\varepsilon_{2}(u)}}
\quad\mathrm{mod}\ \Pi_{0}
\]
where

\[
\diagram{
\begin{tikzcd}[ampersand replacement=\&,column sep=2em]
\operatorname{sg}:\&
\mathcal{G}\arrow[r,"\mathrm{can}"]\&
\mathcal{G}_{3}\arrow[r,"\operatorname{sg}"]\&
\mu.
\end{tikzcd}
}
\]


This shows \deleted{\unsure{that we}} that

\uline{Lemma}
\[
\left\{
\begin{aligned}
&\text{If }x\in\mathcal{G}^{+},\ u\in\mathcal{M}_{0}^{!},
  \quad\text{then }
  \deleted{xux^{-1}\in\mathcal{M}_{0}^{!}}
  \ \text{equivalent conditions}\\
&a)\quad xux^{-1}\in\mathcal{M}_{0}^{!}\\
&b)\quad xux^{-1}\equiv u\quad(\Pi_{0})
       \quad\text{i.e.}\quad [x,u]\in\Pi_{0}\\
&c)\quad .\operatorname{sg}(x)^{
       \inserted{\unsure{\xi_{2}(u)|\varepsilon_{2}(u)}}}
       \deleted{\unsure{\xi_{2}(u)|\varepsilon_{2}(u)}}=1
       \quad\text{i.e.}\quad
       \operatorname{sg}(x)=1\quad\text{\uline{or}}\quad
       \varepsilon_{2}(u)=1
\end{aligned}
\right.
\]

\margin{638}

\scanpage{640}

\deleted{For $x=\rho$, this shows that \unsure{the|the}
\deleted{\unsure{condition}}\inserted{\unsure{first group}}
\unsure{is}}

\uline{Corollary 1} Let $x\in\mathcal{G}^{+\wedge}$.
For $x$ to normalize $\mathcal{M}_{0}^{!}$,
it is necessary and sufficient that we have $\operatorname{sg}(x)=+1$.

\uline{Corollary 2}. The subgroup
$\mathcal{M}_{0}^{\prime\prime!}
=\mathcal{M}''(0)\mathbin{.}\Pi_{0}$
is \deleted{invariant in} normalized by
$\mathcal{G}^{+\wedge}$.

This shows in particular that $\rho$ normalizes
$\mathcal{M}_{0}^{!}$, hence also
$\inserted{\mathcal{M}'(-\bar j)=}
Z_{\rho,0}^{!}=Z_{\rho}\cap\mathcal{M}_{0}^{!}$,
which is therefore invariant in $Z_{\rho}$.
The same therefore holds for its intersection
with $\mathcal{M}^{0}(-\bar j)\times\mu$, i.e.
$\mathcal{M}^{0}(-\bar j)$.
(Or again; this group is equal to
$Z_{\rho}\cap\mathcal{M}_{0}^{\prime\prime!}$,
and $\mathcal{M}_{0}^{\prime\prime!}$ is normalized by $\rho\ldots$).
\uline{Thus the four groups} $(39)$
\uline{are normal in $Z_{\rho}$.} )

To see whether they are normal in
$\unsure{\mathcal{N}_{\rho}^{*}}
=\mu_{\tau'}\mathbin{.}Z_{\rho}$,
it suffices to see whether they are normalized by $\tau'$.
This is clear for $Z_{\rho}^{!}=Z_{\rho}\cap\mathcal{M}^{!}$
-- whence the fact that
$\mathcal{N}_{\rho}^{?}=\mu_{\tau'}\mathbin{.}Z_{\rho}^{!}$
is a subgroup of $\mathcal{N}_{\rho}^{*}$.
To
\deleted{$\mathcal{M}^{0}(-\bar j)=
Z_{\rho}\cap\mathcal{M}_{0}^{\prime\prime!}$, it}
\inserted{see \deleted{\unsure{whether the}} which of the}
three other subgroups
\inserted{considered}
last are normalized, it would therefore be necessary to make explicit
the action of $\tau'$ \deleted{on them.}
on $Z_{\rho}^{!}=\mathcal{M}'(-\bar j)\times\mu$.

\margin{639}

\scanpage{641}
\mspage{623}

Let us note that $\tau\in\mathcal{M}(0)\subset\mathcal{M}_{0}^{!}$
normalizes $\mathcal{M}_{0}^{!}$, so
\deleted{\unsure{one|therefore}}\inserted{it} operates
\deleted{$\unsure{Z}$} on
$\mathcal{M}^{!}
=\underbrace{\mathcal{M}_{0}^{!}}_{\simeq\widetilde{\Gamma}_{\mathbb{Q}}\mathbin{.}\Pi_{0}}
\times\mu$
by the product operation of conjugation by $\tau$ in
$\mathcal{M}_{0}^{!}$, and the identity in $\mu$.
On the other hand we have $\tau'=\tau\sigma$, we have seen that
\[
 (40)\qquad
 \sigma(x)\equiv\varepsilon_{2}(x)x\qquad(\Pi_{0})
 \qquad x\in\mathcal{M}
\]
hence
\[
 (41)\qquad
 \tau'(x)\equiv\varepsilon_{2}(x)\tau(x)\qquad(\Pi_{0})
 \qquad x\in\mathcal{M}
\]
so it follows
\[
 (42)\quad
 \left\{
 \begin{aligned}
 &\text{if }x\in\mathcal{M}_{0}^{!},\quad
   \text{then }\tau'(x)\equiv x\quad(\Pi_{0})
   \inserted{\Longleftrightarrow\tau'(x)\in\mathcal{M}_{0}^{!}}
   \Longleftrightarrow\varepsilon_{2}(x)=1\\
 &\text{in particular, }\mathcal{M}_{0}^{\prime\prime!}
   \text{ is normalized by }\tau
 \end{aligned}
 \right.
\]
Thus \deleted{if $x\in$} applying this to the elements of
$Z_{\rho,0}^{!}=\mathcal{M}'(-\bar j)$ in
$Z_{\rho}^{!}=\mathcal{M}'(-\bar j)\times\mu$,
we find
\[
 (43)\quad
 \left\{
 \begin{aligned}
 &\text{Let }x\in Z_{\rho,0}^{!}=\mathcal{M}'(-\bar j),
   \text{ then }\deleted{\unsure{\tau}}\\
 &\tau'(x)\in Z_{\rho,0}^{!}
   \Longleftrightarrow\varepsilon_{2}(x)=1
 \end{aligned}
 \right.
\]

\uline{Corollary}\quad
$Z_{\rho,0}^{!}\inserted{=\mathcal{M}'(-\bar j)}$
is not normalized by $\tau'$, hence is not invariant in
$\unsure{\mathcal{N}_{\rho}^{*}}$,
but \deleted{the}
$\mathcal{M}^{0}(-\bar j)=Z_{\rho}\cap\mathcal{M}_{0}^{\prime\prime!}$
is normalized by $\tau'$, hence also its product with $\mu$.

Thus the only one of the four subgroups $(39)$ of $Z_{\rho}$
which is not normal in $\mathcal{N}_{\rho}$ is
$\mathcal{M}'(-\bar j)=Z_{\rho,0}^{!}$.
The others give

\margin{640}

\scanpage{642}

rise, by $\frac12$ direct product with $\mu_{\tau'}$,
to \unsure{subgroups|subgroups} of $\mathcal{N}_{\rho}$,
invariant in $\mathcal{N}_{\rho}$.
The smallest of all these subgroups is
$\mathcal{M}^{0}(-\bar j)$, which is of index $24$ in
$\unsure{\mathcal{N}_{\rho}^{*}}$.
The quotient group \deleted{\unsure{is a}}
\[
 \deleted{\unsure{\text{group}}}\qquad
 \mathcal{N}_{\rho}^{*}/\mathcal{M}^{0}(-\bar j)
\]
is therefore a group of order $24$, which moreover receives the group
$\mathbb{D}_{\rho}
\inserted{=\mu_{\tau'}\mathbin{.}L_{\rho}}
\subset\mathcal{N}_{\rho}^{*}$,
normalizer of $L_{\rho}$ in
$\unsure{\operatorname{GL}}(2,\mathbb{Z})$,
a group isomorphic to $D_{6}$, and of order $12$.
The kernel of

\[
\diagram{
 \begin{tikzcd}[ampersand replacement=\&]
 (44)\quad\mathbb{D}_{\rho}
 \arrow[r,hook] \&
 \mathcal{N}_{\rho}^{*}/
 \mathcal{M}^{0}(\bar j)
 \end{tikzcd}
}
\]

is obviously contained in $L_{\rho}=\mathbb{D}_{\rho}\cap Z_{\rho}$
and it is reduced to $1$ for
since we even have
\[
 L_{\rho}\cap Z_{\rho,0}^{!}
 \subset L_{\rho}\cap\mathcal{M}_{0}^{!}
 =L_{\rho}\cap\Pi_{0}=\{1\}.
\]
Thus $\mathbb{D}_{\rho}\inserted{\text{(of order }12\text{)}}$
appears as a subgroup of index $2$ of
$\mathcal{N}_{\rho}^{*}/\mathcal{M}^{0}(-\bar j)$,
just as it is a subgroup of index $2$ of
$\mathcal{G}/\Pi_{0}$. It would be tempting

\margin{641}

\scanpage{643}
\mspage{624}

then to establish an isomorphism
\[
 \mathcal{N}_{\rho}^{*}/\mathcal{M}^{0}(-\bar j)
 \overset{?}{\simeq}\mathcal{G}^{\wedge}/\Pi_{0}
\]
compatible with these embeddings of $\mathbb{D}_{\rho}$.
Let us try it \unsure{this way|also}:

\uline{Proposition}\quad Let us consider the composite homomorphism

\[
\diagram{
 \begin{tikzcd}[ampersand replacement=\&,column sep=large,row sep=large]
 (45)\quad\mathcal{M}\arrow[r,"\Theta_{\mathcal{M}}"]\&
 \operatorname{GL}(2,\widehat{\mathbb{Z}})
 \arrow[r]\&
 \operatorname{GL}(2,\mathbb{Z}/4\mathbb{Z})
 \arrow[r,"\substack{\text{epimorphism of}\\\text{degree }4}"]\&
 \begin{gathered}
 \mathcal{G}^{\unsure{*|\wedge}}/\Pi_{0}\\
 \simeq\mathcal{G}^{\wedge}/\Pi_{0}^{\wedge}
 \end{gathered}\\
 \&\&\underbrace{\operatorname{GL}(2,\mathbb{Z})}_{\mathcal{G}}
 /\mathcal{G}[4]\arrow[u,"\sim"']\&
 \end{tikzcd}
}
\]

\deleted{and the induced homomorphism}
which induces on $\mathcal{G}^{\wedge}\subset\mathcal{M}$
the canonical homomorphism
$\mathcal{G}^{\wedge}\longrightarrow
\mathcal{G}^{\wedge}/\Pi_{0}^{\wedge}$.
Then the homomorphism \deleted{composite}\inserted{induced}

\[
\diagram{
 \begin{tikzcd}[ampersand replacement=\&]
 \mathcal{N}_{\rho}^{*}\arrow[r]\&\mathcal{G}/\Pi_{0}
 \end{tikzcd}
}
\]

has as kernel $\mathcal{M}^{0}(-\bar j)$, and induces by passage
to the quotient an isomorphism

\[
\diagram{
 \begin{tikzcd}[ampersand replacement=\&]
 (46)\quad\mathcal{N}_{\rho}^{*}/\mathcal{M}^{0}(-\bar j)
 \arrow[r,"\sim"]\&\mathcal{G}/\Pi_{0}
 \end{tikzcd}
}
\]


Proof. The elements of $\mathcal{M}^{0}(-\bar j)$
are the elements of the form
\[
 \left\{
 v=\alpha^{-1}u\ \middle|\
 \begin{aligned}
 &u\in\mathcal{M}(0),\quad\alpha\in\Pi_{0}\\
 &u(\rho)=\alpha(\rho)
   \qquad(\text{thus }\varepsilon_{3}(u)=1)\\
 &\varepsilon_{2}(u)=1
 \end{aligned}
 \right\}
\]
they are in $1$-$1$ correspondence with the elements of
\[
 \mathcal{M}^{0}(0)
 =\{u\in\mathcal{M}(0)\mid\varepsilon_{2}(u)=\varepsilon_{3}(u)=1\}.
\]
The image \unsure{of|of the}

\margin{642}
